% !TeX root = ../supplementary.tex

\section{Supplementary A: Extended Method Derivations}
\label{supp:extended_method}

This section provides the technical details that are intentionally abbreviated in the main paper. The goal is to keep the main method compact while still reporting the complete analytical forms used in implementation.

\subsection{ADMK Adaptive Thresholding and Morphological Refinement}
ADMK combines color-space discrepancy analysis and region-level filtering to improve supervision reliability. After obtaining the smoothed discrepancy map $\widetilde{\Delta}$, the threshold is computed per image:
\begin{equation}
\tau = \mu_{\widetilde{\Delta}} + \alpha\,\sigma_{\widetilde{\Delta}},
\label{eq:supp_admk_tau}
\end{equation}
where $\mu_{\widetilde{\Delta}}$ and $\sigma_{\widetilde{\Delta}}$ are the mean and standard deviation of $\widetilde{\Delta}$. This adaptive rule avoids over-masking in mild scenes and under-masking in heavily distorted scenes.

To reduce isolated noise and maintain coherent regions, binary masks are refined by opening and closing:
\begin{equation}
\mathcal{M} = \bigl((\mathcal{M}_0 \ominus \mathcal{S}) \oplus \mathcal{S}\bigr) \oplus \mathcal{S} \ominus \mathcal{S},
\label{eq:supp_admk_morph}
\end{equation}
where $\mathcal{S}$ is the structuring element. The refined mask is used both in loss weighting and in condition construction to ensure consistent filtering across modules.

\subsection{DBLE Details for Background Light and Transmission}
DBLE estimates physically interpretable guidance maps by decomposing slowly varying illumination and locally varying attenuation. Guided filtering is first used to estimate background light. For each local window $\omega_k$, a linear model is defined as
\begin{equation}
q_i = a_k g_i + b_k,\quad i\in\omega_k,
\end{equation}
and optimized by
\begin{equation}
\min_{a_k,b_k}\sum_{i\in\omega_k}(a_k g_i + b_k - p_i)^2 + \epsilon a_k^2.
\label{eq:supp_guided_obj}
\end{equation}
The closed-form solution is
\begin{equation}
a_k = \frac{\mathrm{cov}_{\omega_k}(g,p)}{\mathrm{var}_{\omega_k}(g)+\epsilon},\qquad
b_k = \mu_{\omega_k}(p) - a_k\mu_{\omega_k}(g).
\label{eq:supp_guided_ab}
\end{equation}
Applying this formulation with a large radius gives
\begin{equation}
B^c = \mathrm{GF}(I_{raw}^c; g, r_B, \epsilon_B),\qquad c\in\{r,g,b\}.
\label{eq:supp_dble_b}
\end{equation}

Transmission is then estimated from a normalized dark response:
\begin{equation}
D(\mathbf{x}) = \min_{\mathbf{y}\in\Omega(\mathbf{x})}\ \min_{c\in\{r,g,b\}} \frac{I_{raw}^{c}(\mathbf{y})}{B^{c}(\mathbf{y})+\epsilon},
\end{equation}
\begin{equation}
T_0(\mathbf{x}) = 1 - \omega D(\mathbf{x}),
\label{eq:supp_dble_t0}
\end{equation}
followed by edge-aware refinement:
\begin{equation}
T(\mathbf{x}) = \mathrm{GF}(T_0; g, r_T, \epsilon_T)(\mathbf{x}).
\label{eq:supp_dble_t}
\end{equation}
To model channel-dependent attenuation, a depth proxy is introduced:
\begin{equation}
\hat d(\mathbf{x}) = -\frac{\log(T(\mathbf{x})+\epsilon)}{\bar\beta},\qquad
T^c(\mathbf{x}) = \exp\bigl(-\beta_c\hat d(\mathbf{x})\bigr),
\label{eq:supp_dble_rgb_t}
\end{equation}
which produces a physically interpretable channel-wise transmission map.

A coarse restoration can be written as
\begin{equation}
\widetilde{J}^{c}(\mathbf{x}) = \frac{I_{raw}^{c}(\mathbf{x}) - B^{c}(\mathbf{x})(1-T^{c}(\mathbf{x}))}{\max(T^{c}(\mathbf{x}),T_{min})}.
\label{eq:supp_coarse_restore}
\end{equation}
This image is not treated as the final output. It serves as auxiliary guidance while the rectified flow model performs the final high-fidelity enhancement.

\subsection{Rectified Flow Implementation Details}
The conditional transport process satisfies the continuity equation:
\begin{equation}
\partial_t p_t(\mathbf{x}\,|\,C) + \nabla_{\mathbf{x}}\cdot\bigl(p_t(\mathbf{x}\,|\,C)\,v_\theta(\mathbf{x},C,t)\bigr)=0.
\label{eq:supp_continuity}
\end{equation}
This expression highlights that probability mass is moved by deterministic flow under condition guidance. In addition to Euler integration, a second-order Heun step is available:
\begin{align}
\widetilde{X}_{t_{i+1}} &= X_{t_i} + v_\theta(X_{t_i},C,t_i)\Delta t,\\
X_{t_{i+1}} &= X_{t_i} + \frac{\Delta t}{2}\Big(v_\theta(X_{t_i},C,t_i)+v_\theta(\widetilde{X}_{t_{i+1}},C,t_{i+1})\Big).
\label{eq:supp_heun}
\end{align}
This update reduces numerical truncation error and stabilizes multi-step inference in scenes with strong color attenuation gradients.

Time conditioning uses sinusoidal embedding:
\begin{equation}
\gamma(t)=\big[\sin(\omega_1 t),\cos(\omega_1 t),\dots,\sin(\omega_L t),\cos(\omega_L t)\big],\quad \omega_\ell=\omega_0 2^{\ell-1}.
\label{eq:supp_time_embed}
\end{equation}
The embedding is injected into residual blocks to preserve temporal awareness across scales.

\subsection{Reconstruction Loss Formulation}
Given interpolation state $X_t$, the reconstructed target is
\begin{equation}
\widehat{X}_1(t)=X_t+(1-t)\,v_\theta(X_t,C,t).
\label{eq:supp_x1hat}
\end{equation}
This form enables image-domain supervision without full trajectory rollout at every iteration.

Structural consistency uses SSIM~\cite{r31}:
\begin{equation}
\mathrm{SSIM}(x,y)=\frac{(2\mu_x\mu_y+c_1)(2\sigma_{xy}+c_2)}{(\mu_x^2+\mu_y^2+c_1)(\sigma_x^2+\sigma_y^2+c_2)},
\label{eq:supp_ssim}
\end{equation}
and perceptual consistency uses VGG-19 features~\cite{r32}:
\begin{equation}
\mathcal{L}_{per}=\sum_{\ell\in\mathcal{S}}\frac{1}{N_\ell}\left\|\varphi_\ell(\widehat{X}_1(t)) - \varphi_\ell(X_1)\right\|_1.
\label{eq:supp_per}
\end{equation}
LAB color consistency is defined by
\begin{equation}
\mathcal{L}_{color}=\frac{\left\|\mathcal{M}\odot\big(\mathrm{ab}(\Phi(\widehat{X}_1(t))) - \mathrm{ab}(\Phi(X_1))\big)\right\|_1}{\|\mathcal{M}\|_1+\varepsilon}.
\label{eq:supp_lcolor}
\end{equation}
The three terms complement the flow loss by constraining structure, perceptual realism, and chromatic fidelity simultaneously.

\section{Supplementary B: Additional Quantitative Visualizations}
\label{supp:additional_quant}

This section reports additional visual summaries that are moved from the main paper to keep the primary results section focused. The figures below present a compact view of perceptual behavior, ablation trends, and broader cross-scene restoration performance.

Figure~\ref{fig:supp_perceptual_bars} compares perceptual metrics across representative methods. The bar distribution is consistent with the table-level conclusion in the main paper and highlights the advantage of RF-MCUIE on LPIPS/FID-oriented visual quality.

\begin{figure}[htbp]
    \centering
    \includegraphics[width=0.92\linewidth]{fig_perceptual_bars.pdf}
    \caption{Expanded perceptual comparison on benchmark datasets.}
    \label{fig:supp_perceptual_bars}
\end{figure}

Figure~\ref{fig:supp_ablation_heatmaps} provides an integrated view of module contribution and step sensitivity. The visualization helps interpret why $N=3$ is selected for testing and how physical conditions plus ADMK improve robustness.

\begin{figure*}[htbp]
    \centering
    \includegraphics[width=0.95\textwidth]{fig_ablation_heatmaps.pdf}
    \caption{Heatmap-based ablation analysis for inference-step sensitivity and module contribution.}
    \label{fig:supp_ablation_heatmaps}
\end{figure*}

To complement the representative qualitative examples shown in the main paper, Figure~\ref{fig:supp_overall_comparison} provides a broader visual comparison across six additional scenes. This figure gives a more complete view of how different methods behave under varied degradations and helps connect the metric-level trends in the main paper with concrete restoration phenomena.

\begin{figure*}[htbp]
    \centering
    \includegraphics[width=0.98\textwidth]{fig_supp_overall_comparison.png}
    \caption{Overall qualitative comparison across six additional underwater scenes. The columns are Raw, ZAP, ICSP, UIE-WD, DM-water, LANet, UIEC2Net, and Ours. The visual trends are consistent with the qualitative conclusions reported in the main paper.}
    \label{fig:supp_overall_comparison}
\end{figure*}

The six rows cover green-dominant color casts, bluish scattering, low-contrast seabed scenes, and examples in which local detail recovery is particularly difficult. ZAP and ICSP often leave under-corrected colors or insufficient contrast restoration when scattering becomes severe. UIE-WD and DM-water improve visibility in some cases, but residual color bias and locally unstable enhancement remain visible in challenging regions. LANet and UIEC2Net recover stronger textures in several samples, yet their global color tone is not equally stable across all scenes. In contrast, RF-MCUIE maintains a more balanced restoration behavior by jointly improving visibility, suppressing residual casts, and preserving structural details with fewer obvious artifacts. These observations are consistent with the quantitative advantages reported in the main paper.

\section{Supplementary C: Extra Training and Optimization Details}
\label{supp:training_details}

RF-MCUIE is implemented in PyTorch 1.8 and trained on one NVIDIA RTX A6000 GPU. Input images are cropped to $256\times256$ and normalized to $[-1,1]$. Optimization uses Adam with $\beta_1=0.9$, $\beta_2=0.999$, learning rate $1\times10^{-4}$, batch size 8, and 1,000,000 iterations.

The training process follows a coarse-to-fine schedule. Stage I optimizes transport alignment through $\mathcal{L}_{flow}$, which improves stability in early iterations. Stage II activates the full objective and gradually increases reconstruction weights according to the schedule in the main paper. This transition avoids abrupt loss-scale jumps and improves convergence smoothness.

At inference time, $n=3$ Euler steps are used by default, balancing restoration quality and latency. When strict runtime constraints are imposed, single-step inference can be used as a lightweight deployment mode with moderate quality degradation.
